\section{Related Work}
\label{sec:related}

\subsection{KG Completion and Link Prediction}

Methods such as TransE~\cite{TODO-TransE}, RotatE~\cite{TODO-RotatE}, and
ComplEx~\cite{TODO-ComplEx} learn entity and relation embeddings to predict missing
edges within a single KG.
These approaches assume a shared schema and are designed to densify an existing graph,
not to generate hypotheses that cross between two graphs with disjoint schemas.
Multi-relational embedding approaches similarly operate within a unified ontological
space and cannot represent the structural boundary between two independently curated
domain knowledge graphs.

\subsection{Multi-KG Alignment}

Entity alignment~\cite{TODO-entity-alignment-survey},
schema alignment~\cite{TODO-schema-alignment}, and ontology
matching~\cite{TODO-ontology-matching} address the problem of mapping entities
between knowledge graphs.
These methods focus on constructing the mapping itself, not on exploiting it to
generate cross-domain hypothesis candidates.
Our \textit{align} operator draws on this literature but is operationalised
specifically to enable path-based cross-domain traversal via the subsequent
\textit{union} and \textit{compose} steps.

\subsection{Biomedical KG Discovery}

SemMedDB~\cite{TODO-SemMedDB} and drug repurposing systems such as
PREDICT~\cite{TODO-PREDICT} use predicate-level associations mined from the
biomedical literature for mechanism discovery and drug repurposing.
These systems operate within a unified biomedical ontology (a single-schema KG)
rather than across two independently curated domain graphs.
Production-scale biomedical KGs such as Hetionet~\cite{TODO-Hetionet} and
PrimeKG~\cite{TODO-PrimeKG} represent a single integrated schema and therefore
do not exhibit the structural cross-domain boundary we study.
Our work is complementary: it studies the behaviour of a multi-operator pipeline
on small-to-medium Wikidata-derived KGs, with the goal of characterising the
structural conditions for cross-domain discovery.

\subsection{Operator-Based KG Reasoning}

Description logics~\cite{TODO-description-logics} and rule-based
reasoners~\cite{TODO-rule-reasoning} derive implicit facts within a schema but
do not define operators for cross-KG path composition.
Our work introduces a graph-algebraic operator vocabulary (align, union, compose,
difference, evaluate) as the compositional primitive for cross-domain discovery.

\subsection{Swanson's ABC Model and Literature-Based Discovery}

The classic ABC discovery model~\cite{swanson1986fish} finds implicit connections
$A \to C$ in two separate literatures via a shared intermediate concept $B$.
The seminal example is the fish-oil / migraine / Raynaud's syndrome connection:
literature~1 linked $A$ (fish oil) to $B$ (blood rheology), while literature~2
linked $B$ (blood rheology) to $C$ (Raynaud's syndrome), yet the connection
$A \to C$ was not explicitly stated in either.
Our \textit{align} operator implements a structural analogue: bridge entities $B$
shared between $G_{\text{bio}}$ and $G_{\text{chem}}$ enable $A \to B \to C$
paths that cross the domain boundary.
Our contribution is to formalise this as a KG operator pipeline, to characterise
the conditions under which the mechanism produces useful output, and to identify
the principal source of semantic failures in the generated candidates.
